\problemname{Tomatoes}
An interesting fact is that unripe tomatoes ripen faster if you place some already ripe tomatoes among them.
In this problem, you shall simulate this process and determine how many tomatoes are ripe after a certain number of days.

Assume that $n$ tomatoes lie in a long row and are numbered from $1$ to $n$.
Three of these tomatoes, numbers $t_1$, $t_2$, and $t_3$, are already ripe when the simulation starts at day $0$.
Each day, the tomatoes lying immediately next to an already ripe tomato ripen.
After day $1$, the neighbors of the three initially ripe tomatoes have ripened; after day $2$, the neighbors of those that ripened during day $1$ have also ripened, and so on.

Write a program that, given the number of tomatoes $n$, the number of days $d$, and the numbers $t_1$, $t_2$, $t_3$, computes how many tomatoes are ripe after $d$ days.

\section*{Input}
The first line of input contains two numbers $n$ ($3 \le n \le 100$) and $d$ ($1 \le d \le 100$).

The second line contains the numbers $t_1$, $t_2$, and $t_3$, all distinct and in the range $1 \dots n$.

\section*{Output}
Output a single number: the number of ripe tomatoes after $d$ days.
